\documentclass[conference,a4paper]{IEEEtran}
\usepackage{cite}
\usepackage{amsmath,amssymb,amsfonts}
\usepackage{graphicx}
\usepackage{textcomp}
\usepackage{xcolor}
\usepackage{hyperref}

\begin{document}

\title{Assignment 1: Research Proposal\\Outline \& Annotated Bibliography}

\author{\IEEEauthorblockN{Daniel Tyukov (5714699)}
\IEEEauthorblockA{Department of Microelectronics\\
Delft University of Technology\\
Delft, The Netherlands\\
d.tyukov@student.tudelft.nl}
\and
\IEEEauthorblockN{Apostolos Tsailas (5506387)}
\IEEEauthorblockA{Department of Microelectronics\\
Delft University of Technology\\
Delft, The Netherlands\\
a.tsailas-1@student.tudelft.nl}
\thanks{This article was written as part of the course EE4C01 Profile Orientation \& Academic Skills of the faculty Electrical Engineering, Mathematics and Computer Science of Delft University of Technology.}
}

\maketitle

% ============================================================
\section{Research Question and Motive}
% ============================================================

\textbf{Main question:} To what extent does backside power delivery improve signal integrity and reduce electromagnetic coupling compared to conventional front-side power distribution in advanced CMOS nodes?

\medskip

\textbf{Motive:} Below 5\,nm, the wires on a chip get so resistive that delivering clean power through the front side becomes a real bottleneck. The power rails and signal lines all share the same crowded metal stack, which causes IR drop issues and unwanted coupling between them. Recently, companies like Intel have started moving the power network to the backside of the wafer. Intel calls their version PowerVia, and imec has been publishing work on buried power rails paired with nano-TSVs. These approaches look promising, but they also bring new problems. Thinning the wafer to a few hundred nanometers hurts thermal spreading, the tiny TSV connections add parasitic resistance, and the fabrication itself can warp the wafer. What we noticed is that there is no paper yet that pulls all of this together and critically weighs the signal integrity gains against these new trade-offs. That is the gap we want to fill with our literature review.

\subsection{Sub-questions}
\begin{enumerate}
    \item[A.] What are front-side power delivery networks (FSPDNs) and what signal integrity problems do they cause at advanced nodes?
    \item[B.] How are backside power delivery networks (BSPDNs) built and what makes them structurally different from FSPDNs?
    \item[C.] What happens to IR drop, supply noise, and voltage droop when you move the power rails to the backside?
    \item[D.] Which electromagnetic coupling effects between power and signal lines does BSPDN help reduce?
    \item[E.] What new signal integrity risks come with BSPDN, for instance nanoTSV parasitics or self-heating from the thinned substrate?
    \item[F.] What simulation and measurement approaches are used to compare signal integrity between FSPDN and BSPDN?
    \item[G.] How do Intel's PowerVia and imec's buried-power-rail approach compare in terms of signal integrity results?
    \item[H.] What open questions remain, and where should future work focus to fully understand BSPDN signal integrity?
\end{enumerate}

% ============================================================
\section{Writing Plan}
% ============================================================

\subsection{Proposed Structure}

The paper will have five sections. We open with an introduction that sets up the scaling problem and explains why backside power delivery has gained traction (sub-questions A, B, and the main question). A short method section follows, describing our literature search: which databases we used, what keywords, and our selection criteria.

The body will have three sections. First, we cover the limitations of conventional front-side PDN and make the case for why the industry is moving to backside delivery (A, B). Then we discuss the actual signal integrity improvements that have been reported: lower IR drop, less coupling, and how different implementations compare (C, D, G). A third body section covers the flip side: new problems BSPDN introduces, like nanoTSV parasitics, thermal penalties, and wafer warpage, plus how researchers have tried to model these effects (E, F).

In the discussion and conclusion we ask: are the signal integrity improvements worth the extra process steps and the thermal hit? We also lay out what remains unresolved (G, H, main question).

% ============================================================
\section{Annotated Bibliography}
% ============================================================

% Reference 1 (Daniel Tyukov)
\subsection{Reference 1}
W.~Hafez \textit{et al.}, ``Intel PowerVia technology: Backside power delivery for high density and high-performance computing,'' in \textit{Proc. IEEE Symp. VLSI Technol. Circuits}, Kyoto, Japan, Jun.~2023, doi: 10.23919/VLSITechnologyandCir57934.2023.10185208.

This is the most-cited industry paper on backside power delivery, with over 50 citations since 2023. Intel built PowerVia on top of their Intel~4 finFET process: instead of sharing the front-side metals between power and signals, they run power from the back of the wafer directly to the transistor source/drain contacts. The front side then has all 14 metal layers free for signal routing.

\textbf{Motive:} At advanced nodes, designers are forced to give up signal routing tracks for power delivery. This trade-off caps both the density and the performance of the chip.

\textbf{Objective:} Demonstrate a working backside power delivery technology on Intel~4 and measure how much it actually helps in terms of cell utilization, voltage droop, and clock speed.

\textbf{Conclusion:} They report more than 90\% cell utilization on a fabricated E-core, a 30\%+ reduction in voltage droop from package to transistor, and a 6\% clock speed boost over the same design without PowerVia. Transistor reliability metrics (TDDB, BTI, HCI) and manufacturing yield were not affected by the backside process.

\textbf{Argumentation:}
\begin{enumerate}
    \item The PowerVia connection sits coplanar with the source/drain contacts, so power reaches the transistor through a much shorter, lower-resistance path than in a conventional front-side PDN.
    \item They measured a 5$\times$ simulated IR drop reduction from bump to transistor. On-chip monitors confirmed over 30\% droop improvement at the package level.
    \item Yield tracks the baseline Intel~4 process closely, which suggests the technology is ready for volume production.
\end{enumerate}

\textbf{Help:} We will use this paper heavily for sub-questions B, C, and G. It gives us hard numbers on what backside delivery can achieve in practice, and other BSPDN approaches can be measured against it.

\bigskip

% Reference 2 (Apostolos Tsailas)
\subsection{Reference 2}
D.~Prasad, S.~S.~Teja~Nibhanupudi, S.~Das, O.~Zografos, B.~Chehab, and S.~Sarkar, ``Buried power rails and back-side power grids: Arm CPU power delivery network design beyond 5nm,'' in \textit{Proc. IEEE Int. Electron Devices Meeting (IEDM)}, San Francisco, CA, USA, Dec.~2019, doi: 10.1109/IEDM19573.2019.8993617.

This is one of the earliest and most-cited papers (69 citations) that systematically compares three PDN architectures: the conventional front-side approach, front-side with buried power rails, and full backside delivery with buried rails. The authors used the Arm Cortex-A53 CPU as their benchmark, placed and routed at imec's 3\,nm node, which makes the comparison very concrete.

\textbf{Motive:} At 3\,nm, interconnect resistance has grown to the point where no amount of power grid optimization on the front side can hold IR drop within acceptable margins for a high-performance CPU.

\textbf{Objective:} Run a complete PDN design study covering both DC (resistive drop) and AC (resonance behavior) across three architectures to see where backside delivery really helps and where it might cause new issues.

\textbf{Conclusion:} Backside delivery with buried rails slashes IR drop by about 7$\times$ compared to conventional front-side PDN, bringing it down to roughly 1\% of V\textsubscript{DD}. But there is a catch: the very low PDN resistance reduces damping, so the chip-package system can ring at resonance. The authors show that on-chip decoupling capacitors mitigate this, especially in multi-core designs where idle cores act as decap.

\textbf{Argumentation:}
\begin{enumerate}
    \item Even the most optimized front-side PDN at 3\,nm still exceeds the 10\% V\textsubscript{DD} IR drop budget, which is a strong argument that front-side delivery is hitting a wall.
    \item With tight $\mu$TSV pitch, backside delivery gets IR drop to around 1\% of V\textsubscript{DD}. At the same time, it frees every front-side metal layer for signals. No more fighting for routing resources.
    \item Their AC analysis matters most for our purposes. It shows a trade-off that most BSPDN papers ignore: lower resistance helps DC drop but hurts damping at resonance.
\end{enumerate}

\textbf{Help:} This source is central to sub-questions A, C, and F. It gives us the methodology template for comparing PDN architectures and -- unusually for BSPDN work -- also looks at AC resonance behavior. That resonance trade-off is something most BSPDN papers skip, and our review should give it proper attention.

\end{document}
